\chapter*{如何使用本书}
\addcontentsline{toc}{chapter}{如何使用本书}
\markboth{如何使用本书}{如何使用本书}

\section*{这是一种什么样的书？}

本书既不是百科全书式教材，也不是私人公式表。它是一组经过结构化整理的本科数学笔记，主要服务于两个用途：
\begin{enumerate}
    \item 作为\textbf{有引导的复习材料}：当读者已经学过某个主题，但希望看见更清楚的联系、证明和例子时使用；
    \item 作为\textbf{参考地图}：当读者需要定位一个定义、定理、技术或应用，并理解它在更大理论中的位置时使用。
\end{enumerate}
一门完整的第一课程仍然需要课堂讲授、习题课和专门教材。本手稿中的练习旨在让阅读变得主动，而不是模拟一整学期的考核训练。

\section*{推荐阅读路线}

\begin{description}[leftmargin=3.3cm,style=nextline]
    \item[标准本科路线] 从第 2 章“数学分析 I”和第 3 章“线性代数”开始；随后按兴趣进入第 4--7 章。第 0--1 章可作为可选参考材料。
    \item[基础路线] 先读第 0--1 章，再进入第 2 章。若数系构造、形式逻辑与公理化集合论本身就是学习目标，这条路线更合适。
    \item[计算与数据路线] 先读第 3 章，再读第 8、9 章；需要优化、逼近、积分和 Fourier 分析时，再回到第 2、5 章。
    \item[结构路线] 阅读第 3、4、6、7 章。这条路线强调向量空间、代数结构、拓扑不变量与全纯结构。
\end{description}

\section*{依赖关系图}

大多数章节在局部上相对自足，但以下依赖关系很有帮助：
\[
\begin{array}{ccl}
\text{数学分析 II} &\longleftarrow& \text{数学分析 I},\\
\text{抽象代数} &\longleftarrow& \text{基础线性代数},\\
\text{拓扑学} &\longleftarrow& \text{集合、函数与初等分析},\\
\text{复分析} &\longleftarrow& \text{数学分析 I 与基础拓扑},\\
\text{概率论} &\longleftarrow& \text{微积分；测度论部分初读时可选读}.
\end{array}
\]
图论几乎可以独立阅读。

\section*{证明与范围约定}

若一个定理后面跟有\textbf{证明}，则该证明应当在本章层次上构成一个完整论证。若某个结果被称为\textbf{高级定理}，则它主要用于提供方向感；完整证明通常需要相当多的额外工具。被省略的论证不应被误认为只是一个初等练习。

只要有助于理解，定义后面会配有例子、非例子或计算解释。应用框展示一个概念如何进入科学、工程或数据分析中的模型；它们并不声称这些应用已经被所展示的数学完全穷尽。

\section*{如何处理练习}

练习大致按照从直接验证到综合运用的顺序排列。一个较好的最低限度训练流程是：
\begin{enumerate}
    \item 凭记忆重构一个证明；
    \item 完成两个计算练习；
    \item 完成一个概念题或反例题；
    \item 用自己的话解释一个应用。
\end{enumerate}
“提示与答案选摘”有意保持简短。它们提供路线或检查点，而不是完整书面解答的替代品。

\begin{mynote}
基础材料在能够澄清后续构造时非常有用，但它不应成为学习微积分、线性代数、图论或概率论的入场费。应当从你的数学问题真正开始的地方进入；只有当某个依赖确实重要时，再向后追溯。
\end{mynote}
